\introduction

Формальные спецификации применяются для описания допустимого поведения сложных
программных систем, в особенности распределённых алгоритмов, для которых
характерны недетерминизм, конкурентное взаимодействие и большое число
возможных сценариев выполнения. Одним из распространённых средств такого
описания является язык TLA+, позволяющий задавать систему в терминах состояний
и переходов, а также проверять свойства модели средствами модельной проверки.

Однако наличие формальной спецификации само по себе не гарантирует её
актуальность по отношению к реальной реализации. В процессе разработки
исходный код системы изменяется, дополняется новыми возможностями и
оптимизируется, тогда как формальная модель нередко обновляется с запозданием
или не обновляется вовсе. В результате возникает риск расхождения между
фактическим поведением программы и поведением, допускаемым спецификацией.
Подобное расхождение снижает ценность формальной модели как инструмента
анализа и верификации.

Одним из возможных способов решения данной проблемы является валидация трасс
исполнения программы относительно формальной спецификации. В этом подходе
выполнение программы наблюдается и фиксируется в виде трассы, содержащей
сведения о происходящих событиях и изменениях логических переменных. Затем эта
трасса интерпретируется в терминах формальной модели, после чего проверяется,
существует ли допустимое поведение спецификации, согласованное с наблюдаемым
исполнением. Такой подход позволяет сравнивать реализацию и спецификацию без
необходимости полного доказательства корректности программы относительно модели.

Актуальность работы определяется необходимостью построения формального и
практически реализуемого метода проверки согласованности между исходным кодом
и TLA+-спецификацией. Особенно важна такая проверка для распределённых систем,
где вручную сопоставить поведение реализации и модели затруднительно из-за
сложности пространства состояний и большого числа возможных межпроцессных
взаимодействий.

Целью работы является разработка методики верификации соответствия трасс
исполнения программы и формальной спецификации на языке TLA+, а также
формализация математической постановки этой задачи и обоснование её решения.

Место прохождения практики – Научно-учебный комплекс «Информатика и системы
управления» МГТУ им. Н.Э. Баумана. Период прохождения практики – с 9 февраля
2026 года по 10 марта 2026 года.

В задачи прохождения практики входило:

\begin{enumerate}
    \item Сформулировать математическую постановку задачи верификации
    соответствия TLA+ спецификации и исходного кода;
    \item Предложить возможные методы решения поставленной задачи.
\end{enumerate}

Практическая значимость работы заключается в том, что предложенный подход
может использоваться для автоматизированного выявления расхождений между
реализацией и формальной моделью, а также для локализации шага трассы, на
котором впервые возникает несоответствие.

Структурно работа состоит из разделов, посвящённых описанию методики
верификации, математической постановке задачи, обоснованию её решения и
реализации полученных результатов средствами TLA+ и TLC.
